% ===================================================================
%  LIMINAIRES — feuillets 1 a 6 : couverture, page de titre, dedicace
%  + KONSTATO de l'Akademio Idista, PREFACO.
%  Aucune de ces pages ne porte de folio imprime, SAUF le feuillet 6
%  (2e page de la PREFACO), qui porte « — 4 — » : c'est une page de
%  continuation, non une page d'ouverture (meme convention que dans
%  les tableaux : voir texto/io/10-tabelo-01.tex, feuillet 7 sans
%  folio / feuillet 8 avec folio 6).
%
%  Corps et interlettrages ESTIMES A L'ŒIL sur le rapport a la hauteur
%  de capitale du corps courant (10,93 pt = 54 px de haut sur le scan
%  brut, feuillet 8), PROVISOIRES comme le reste du releve d'apparat.
% ===================================================================

% --- Feuillet 1 (couverture, non foliote) ---------------------------
%  Encadrement/ornement : aucun filet de bordure autour de la page ;
%  seuls deux filets horizontaux centres marquent le texte (voir
%  ci-dessous). Le premier, sous IDO-AKADEMIANO, est orne de volutes a
%  ses deux extremites (non reproduites, ~99 mm) ; le second, sous le
%  sous-titre, est un FILET DOUBLE (deux traits paralleles rapproches,
%  ~40 mm), rendu ici par deux \VUfilet successifs tres rapproches.
\begin{VUpage}[1]{}
%%K t00-apar-1 apar
\VUsaut{4.0mm}
\VUcentre{16.4pt}{40}{J. GUIGNON}
\VUsaut{2.7mm}
\VUcentre{7.7pt}{140}{IDO-AKADEMIANO}
\VUsaut{2.5mm}
\VUfilet{78mm}
\VUsaut{14.5mm}
\VUcentre{22.5pt}{85}{EXPLIKO=LIBRETO}
\VUsaut{4.7mm}
\VUcentre{8.5pt}{160}{DI}
\VUsaut{5.1mm}
\VUtitre{19.0pt}{0}{La Delmas-Tabeli Helpanta}
\VUsaut{6.3mm}
\VUcentre{10.1pt}{0}{por}
\VUsaut{4.4mm}
\VUcentre{18.1pt}{0}{la praktikal docado}
\VUsaut{2.8mm}
\VUcentre{18.1pt}{0}{di la moderna lingui}
\VUsaut{4.0mm}
\VUcentre{15.4pt}{0}{per la direta metodo e l'imaji}
\VUsaut{9.8mm}
\VUfilet{40mm}
\VUsaut{0.4mm}
\VUfilet{40mm}
\VUsaut{27.7mm}
\VUcentre{10.5pt}{80}{IDO-KONTORO}
\VUsaut{2.8mm}
\VUcentre{9.5pt}{60}{THAON-LES-VOSGES}
\VUsaut{1.6mm}
\VUcentre{9.5pt}{60}{\textsc{Vosges --- France}}
\end{VUpage}

% --- Feuillet 2 (verso de la couverture, non foliote) ---------------
%  Page blanche : aucune encre, seul le grain du papier et le
%  transparence du feuillet 1. Position conventionnelle du faux-titre,
%  mais rien n'y est imprime dans cet exemplaire.
\begin{VUpage}[2]{}
% (page vierge)
\end{VUpage}

% --- Feuillet 3 (page de titre, non foliotee) -----------------------
%  Meme composition que le feuillet 1 (couverture), au mot pres. Porte
%  en bas a droite la marque d'imprimeur « 1. Ido. » (repere de
%  cahier) : non transcrite, ce n'est pas une matiere de lecture.
\begin{VUpage}[3]{}
%%K t00-apar-2 apar
\VUsaut{4.0mm}
\VUcentre{16.4pt}{40}{J. GUIGNON}
\VUsaut{2.7mm}
\VUcentre{7.7pt}{140}{IDO-AKADEMIANO}
\VUsaut{2.5mm}
\VUfilet{78mm}
\VUsaut{14.5mm}
\VUcentre{22.5pt}{85}{EXPLIKO=LIBRETO}
\VUsaut{4.7mm}
\VUcentre{8.5pt}{160}{DI}
\VUsaut{5.1mm}
\VUtitre{19.0pt}{0}{La Delmas-Tabeli Helpanta}
\VUsaut{6.3mm}
\VUcentre{10.1pt}{0}{por}
\VUsaut{4.4mm}
\VUcentre{18.1pt}{0}{la praktikal docado}
\VUsaut{2.8mm}
\VUcentre{18.1pt}{0}{di la moderna lingui}
\VUsaut{4.0mm}
\VUcentre{15.4pt}{0}{per la direta metodo e l'imaji}
\VUsaut{9.8mm}
\VUfilet{40mm}
\VUsaut{0.4mm}
\VUfilet{40mm}
\VUsaut{27.7mm}
\VUcentre{10.5pt}{80}{IDO-KONTORO}
\VUsaut{2.8mm}
\VUcentre{9.5pt}{60}{THAON-LES-VOSGES}
\VUsaut{1.6mm}
\VUcentre{9.5pt}{60}{\textsc{Vosges --- France}}
\end{VUpage}

% --- Feuillet 4 (dedicace + KONSTATO, non foliote) ------------------
%  Haut de page : dedicace en italique a la marquize L. de Beaufront.
%  Bas de page : KONSTATO de l'Akademio Idista, presentee dans un
%  encadrement rectangulaire (filet de bordure non reproduit), suivie
%  de la mention legale de depot.
\begin{VUpage}[4]{}
%%K t00-apar-3 apar
\VUsaut{3.0mm}
\VUcentre{13.5pt}{0}{\textit{A}}
\VUsaut{3.2mm}
\VUcentre{13.5pt}{0}{\textit{Sioro Markezo L. de Beaufront}}
\VUsaut{2.8mm}
\VUcentre{13.5pt}{0}{\textit{Ni dedikas la yena libreto}}
\VUsaut{2.1mm}
\VUcentre{13.5pt}{0}{\textit{Kom modest expresuro di gratitudo,}}
\VUsaut{1.5mm}
\VUcentre{13.5pt}{0}{\textit{Okazione di lua sepadekesma aniversario.}}
\VUsaut{2.4mm}
\VUtitre{15.5pt}{0}{3 oktobro 1925, la Vogesiana Idisti.}
\VUsaut{12.1mm}
% Debut de l'encadrement rectangulaire (non reproduit).
\VUtitre{8.3pt}{40}{KONSTATO}
\VUsaut{2.0mm}
% L'encart KONSTATO est compose dans un corps plus petit que le corps
% courant (le mot KONSTATO lui-meme mesure moins haut que la capitale
% du corps courant sur le scan) : estime ici a 8.2pt.
{\fontsize{8.2pt}{9.8pt}\selectfont
\noindent L'Akademio Idista konstatas, ke la linguala formo pri\cc
zentita en ica " \VUgras{Ido-Expliko-Libreto} " esas la oficala\nl
formo di la linguo internaciona Ido, quala ol rezultis de la\nl
labori di la Linguala komitato di la Delegitaro por\nl
l'adopto di Linguo Internaciona, di lua konstanta komi\cc
sitaro e la decidi dil Akademio Idista.

\VUsaut{2.0mm}
{\raggedleft\textit{La prezidanto,}\par}
\noindent 8 février 1926.\hfill F. SCHNEEBERGER.\par}
\VUsaut{6.0mm}
% Fin de l'encadrement rectangulaire (non reproduit).
{\fontsize{8.0pt}{9.6pt}\selectfont
\VUblancAlinea
\VUgras{Le dépôt du présent ouvrage a été effectué conformément aux lois}\nl
\VUgras{françaises et aux conventions internationales.}

\VUblancAlinea
\VUgras{Tous droits réservés. --- Reproduction interdite.}}
\end{VUpage}

% --- Feuillet 5 (ouverture de la PREFACO, non foliote) --------------
\begin{VUpage}[5]{}
%%K t00-tit-1 sub
\VUsaut{3.0mm}
\VUtitre{18.8pt}{60}{PREFACO}
\VUsaut{2.0mm}
\VUfilet{14.5mm}
\VUsaut{6.0mm}

%%K t00-01-1 p
On postulas de me prefacizar ta libreto !

%%K t00-02-1 p
\VUblancAlinea
Ka prefacizar esas utila ? Kad on lektas la\nl
prefaci ?

%%K t00-03-1 p
\VUblancAlinea
Se me esus certa pri tala lekto, me skribus\nl
kelka linei por indikar la skopo di ca libreto.

%%K t00-04-1 p
\VUblancAlinea
Nu ! me supozez ke on lektos ta prefaco ! Yen\nl
lo dicenda.

%%K t00-05-1 p
\VUblancAlinea
La skopo di ca libreto esas satisfacor l'ex\cc
pekto dil Idistaro.

%%K t00-06-1 p
\VUblancAlinea
Omnaloke ed omnadie nia samideani repe\cc
tas : « Dum mult yari me lektis Ido-revui, e me\nl
tre bone komprenis e komprenas oli, ma kande\nl
me renkontras samideano, me ne sucesas\nl
fluante konversar pri irga temo ! »

%%K t00-07-1 p
\VUblancAlinea
Ni inquestis ed agnoskis ta nehabileso; e\nl
yen, segun nia judiko la « pro quo » :

%%K t00-08-1 p
\VUblancAlinea
Irga revuo ofte traktas temi pri cienco pura,\nl
pri etiko, pri arto o pri literaturo, ed ordinare\nl
la sama vorti sempre rilektesas e riaparas tra\nl
olua pagini. --- On povas lektar revui e multa\nl
libri, nultempe on trovos en oli la konver\cc
sala vorti komunuza ed omnadia, nultempe on\nl
renkontros ula vorti quale bracho, bretelo, kal\cc
drono, korkotirilo, e. c., e. c...

%%K t00-09-1 p
\VUblancAlinea
Tamen Ido devas esar ne nur linguo cien\cc
cala, artala, etikala, ma anke ol devas esar\nl
praktikala linguo.

%%K t00-10-1 p
\VUblancAlinea
Ni opinionis ke mikra libro esus utila al\nl
Idistaro : libreto en qua esus maxim multa\nl
vorti komunuza en tre poka pagini !

%%K t00-11-1 p
\VUblancAlinea
Taspeca libreton on darfus lektar plurfoye\nl
singlayare e pro ta lektado on merkus la vorti\nl
di la ordinara vivo.
\end{VUpage}

% --- Feuillet 6 (suite de la PREFACO, folio 4) -----------------------
\begin{VUpage}[6]{4}
%%K t00-12-1 p
Ta libreto dezirinda esas la yena.

%%K t00-13-1 p
\VUblancAlinea
Irgu qua uzabos lu kom Ido- « vademecum »,\nl
rapide sucesos fluante parolar nia helpo-linguo\nl
internaciona.

%%K t00-14-1 p
\VUblancAlinea
Ni dankas l'editero sioro Delmas qua yuri\cc
zis ni uzar lua praktikala libreto, precipue pro\nl
ke ol esas facile uzebla en irga linguo nacio\cc
nala. Fakte, de nun (po la sama preco kam ta\nl
dil Ido-libreto) on povas komendar de la Ido\cc
Kontoro en Thaon-les-Vosges (Vosges) la ko\cc
respondanta libreti Angla, Flandrana, Franca,\nl
Germana, Hispana, Italiana, Rusa. Uzar ica\nl
tradukuro esus marveloza moyeno praktikala,\nl
nam pos lektir e rilektir la Ido-libreto, kom\cc
parante kun sua nacionala linguo; pos sucesir\nl
komprenar omno, singlu obtenus mem plu\nl
bona rezultajo uzante la libreto nacionala exer\cc
cante su per tradukar a Ido la vorti parlek\cc
tita.

%%K t00-15-1 p
\VUblancAlinea
La numeri relatas kun la belega tabeli qui\nl
editesis segun duopla modelo : murala tabeli\nl
e kayera tabeli. Ta tabeli esas utila precipue\nl
en la Ido-kursi publika, segun la direta me\cc
todo, t. e. per praktikala konversado pri la\nl
temi sugestita da la imaji e diversa ceni.

%%K t00-16-1 p
\VUblancAlinea
Nun ni diletas deklarar ke ni prizentas ica\nl
verketo kom ora buketo a sioro markeza L. de\nl
Beaufront, okazione di lua sepa dek yari. Ol\nl
proklamos koram l'Idistaro : « Konstatez\nl
quante praktikala, logikoza e harmonioza esas\nl
ta linguo helpanta di qua il esis ed esas la patro\nl
tante experta e devota. »

%%K t00-17-1 p
\VUblancAlinea
Dezirez, ni omna dezirez ke il povez longa\cc
tempe parlaborar ta splendid entraprezo qua\nl
vere esas un de la maxim bela inter la homala\nl
solvuri.\hfill J. GUIGNON.
\end{VUpage}
